\documentclass[11pt,a4paper]{article}
\usepackage[T1]{fontenc}
\usepackage[utf8]{inputenc}
\usepackage{lmodern}
\usepackage{microtype}
\usepackage{geometry}
\usepackage{amsmath,amssymb}
\usepackage{booktabs,longtable,tabularx,array}
\usepackage{enumitem}
\usepackage[dvipsnames,table]{xcolor}
\usepackage{hyperref}
\usepackage{xurl}

\geometry{margin=2.25cm}
\hypersetup{
  colorlinks=true,
  linkcolor=MidnightBlue,
  urlcolor=MidnightBlue,
  pdftitle={Relatório de correções e reconciliação das orientações},
  pdfauthor={Revisão técnica assistida},
  pdfsubject={Biblioteca determinística de segmentos afins em células diádicas},
  pdfkeywords={auditoria, correções, aproximação, Besov, complexidade de bits}
}
\setlist[itemize]{leftmargin=1.55em,itemsep=0.25em,topsep=0.35em}
\setlist[enumerate]{leftmargin=1.8em,itemsep=0.3em,topsep=0.35em}
\setlength{\parindent}{0pt}
\setlength{\parskip}{0.55em}
\renewcommand{\arraystretch}{1.16}
\newcolumntype{L}[1]{>{\raggedright\arraybackslash}p{#1}}
\newcommand{\feito}{\textcolor{ForestGreen}{\textbf{Concluído}}}
\newcommand{\parcial}{\textcolor{BurntOrange}{\textbf{Materializado / pendente externo}}}
\newcommand{\externo}{\textcolor{BrickRed}{\textbf{Depende do autor ou de serviço externo}}}
\newcommand{\Library}{\mathcal L}
\newcommand{\LibraryMarked}{\mathcal L^{\square}}

\begin{document}
\pagenumbering{roman}
\hypersetup{pageanchor=false}

\begin{titlepage}
\centering
\vspace*{2.4cm}
{\LARGE\bfseries Relatório de correções e reconciliação das orientações\par}
\vspace{0.65cm}
{\Large Biblioteca determinística de segmentos afins em células diádicas\par}
\vspace{0.35cm}
{\large endereçamento, seleção exata e teoria direta--inversa\par}
\vspace{1.4cm}
{\footnotesize
\begin{tabular}{@{}rl@{}}
Documento de autoridade inicial: & \texttt{parecer.pdf}\\
Revisão posterior: & \texttt{plano.pdf} (versão 2)\\
Pacote corrigido: & \path{Biblioteca_Deterministica_Artigo_Unificado_CORRIGIDO}\\
Data de fechamento: & 8 de agosto de 2026\\
\end{tabular}
}
\vfill
{\small Este relatório registra decisões editoriais verificáveis, alterações matemáticas,
substituições textuais, arquivos afetados, testes e pendências. Não registra nem pretende
reconstruir raciocínio privado; apresenta as justificativas técnicas e as evidências
necessárias para auditoria.\par}
\end{titlepage}
\hypersetup{pageanchor=true}

\tableofcontents
\clearpage
\pagenumbering{arabic}

\section{Resultado entregue}

Foi produzida uma versão corrigida e recompilável do documento-sistema. O PDF principal
tem 77 páginas A4 e foi gerado a partir de um único mestre ativo. As correções materiais
incluem: reposicionamento do funcional de realização; uniformidade quantitativa sobre bolas
de Besov; transferência da testemunha de agudeza para a biblioteca original; fechamento
de uma cota de complexidade de bits; regra exata de desempate; reclassificação da camada de
programação dinâmica como teoria auxiliar de exatidão; ampliação da introdução; integração
dos artefatos antes órfãos; novas referências; declarações editoriais; e endurecimento das
auditorias automatizadas.

O resultado deve ser entendido em três camadas:

\begin{enumerate}
  \item \textbf{Documento-sistema corrigido.} Esta é a versão integral preservada pelo
  guarda-corpo G-6 do plano revisado.
  \item \textbf{Arquitetura editorial congelada.} Os escopos A, B, C e R foram
  materializados em mapas separados no diretório \path{publicacao/}, sem remover nenhum
  teorema do documento unificado.
  \item \textbf{Ações externas não simuladas.} Depósito em arXiv/Zenodo, DOI, ORCID,
  confirmações autorais e submissão a periódico não foram inventados. A tradução integral e
  a extração final de manuscritos independentes continuam assinaladas como etapas
  editoriais, não como fatos consumados.
\end{enumerate}

\section{Base documental e regra de precedência}

\subsection{Entradas controladas}

\begin{longtable}{L{0.29\textwidth}L{0.58\textwidth}}
\toprule
\textbf{Entrada} & \textbf{SHA-256 e função na revisão}\\
\midrule
\endfirsthead
\toprule
\textbf{Entrada} & \textbf{SHA-256 e função na revisão}\\
\midrule
\endhead
Pacote-base auditado &
\nolinkurl{791b73fac4b128e21c999daaad205769420ec8ed443d32670df2025c555258a9}. Fonte
primária dos arquivos corrigidos.\\
\texttt{parecer.pdf} &
\nolinkurl{c9e68529e7f5ed80f591e30566cbbc35ceedbce75aa04a6cb1b58e670e15499d}.
Diagnóstico técnico original: M-1--M-6, B-1--B-3 e E-1--E-8.\\
\texttt{plano.pdf} &
\nolinkurl{25c249d1241e9800e06e0b7c23a8d7235fb8ea236ce7dd7c2896182da176dee2}.
Revisão posterior da estratégia editorial e das dependências do DLPG.\\
Relatório de fortalecimento matemático &
\nolinkurl{39b6f449effb23f9a393d1abba729a0f26b34317bf87949a321b4bacfed77fce}.
Material de apoio para a verificação das formulações matemáticas.\\
Pesquisa de artigos científicos &
\nolinkurl{507f295db2036de465a0158b72a80db92cc07a33ccc11f45d13ad82b826e0be6}.
Material auxiliar para triangulação bibliográfica.\\
Pacote anterior dos itens 1 a 7 &
\nolinkurl{c266a788c1e82062892cacb527caa0572106399243a2575f833c48a1448f29b4}.
Referência histórica de escopo; não substituiu o pacote-base.\\
\bottomrule
\end{longtable}

\subsection{Interpretação reconciliada}

O cabeçalho do plano revisado diz expressamente que ele \emph{substitui o plano da
Seção 6 do parecer}, e o próprio texto declara que nenhum achado matemático foi
invalidado. Dessa redação resultou a seguinte regra, aplicada de forma conservadora:

\begin{enumerate}
  \item os achados M-1--M-6, B-1--B-3 e E-1--E-8 permaneceram obrigatórios;
  \item M-3 foi reforçado: a transferência de $\LibraryMarked$ para $\Library$ deixou de
  ser opcional;
  \item M-4 foi estreitado: a crítica recai sobre as recorrências e sua moldura empírica,
  não sobre o catálogo finito e o desempate da antiga Seção 8.6;
  \item a divisão passou de três artigos para as peças A/B/C, com D (DLPG) dependente de
  B e o residual R destinado preferencialmente a um apêndice de exatidão;
  \item o documento unificado deveria sobreviver como documento-sistema;
  \item J-1--J-3 tornaram obrigatórias explicações internas sobre $\Gamma$ versus
  $\Gamma^\sharp$, inclinação não nula e relevância do refinamento adaptativo.
\end{enumerate}

Assim, o plano novo não foi usado para apagar críticas do parecer: foi usado para mudar
a arquitetura e aumentar o rigor de M-3, preservando o restante.

\section{Matriz ponto a ponto}

\small
\begin{longtable}{L{0.07\textwidth}L{0.36\textwidth}L{0.45\textwidth}}
\toprule
\textbf{ID} & \textbf{Leitura consolidada} & \textbf{Correção e evidência}\\
\midrule
\endfirsthead
\toprule
\textbf{ID} & \textbf{Leitura consolidada} & \textbf{Correção e evidência}\\
\midrule
\endhead
M-1 & O certificado é um funcional de realização; a recíproca é universal e não
caracteriza aproximabilidade específica da biblioteca. & \feito. Citação a
Hristov--Ivanov; Seção 20 reenquadrada; novo lema de desigualdade inversa universal;
observação explícita sobre o que o funcional mede; lado direto com agenda de altura em
primeiro plano.\\
M-2 & O limiar dependente de $f$ precisava ser uniformizado sobre bolas é convertido em
cota quantitativa. & \feito. Nova racionalização por grades alternadas, constante
$A_{K,s}$ independente de $f$, agenda calibrada por $B$ e corolário de imersões
uniformes sobre bolas.\\
M-3 & A agudeza devia valer na biblioteca original $\Library$, não somente em
$\LibraryMarked$. & \feito. Cada folha plana do pente é subdividida e recebe uma tenda
diádica de altura $\delta_j$; as novas inclinações são não nulas; folhas, altura e erro
são controlados. O script de reversibilidade audita a transferência.\\
M-4 & O catálogo finito e o desempate continuam essenciais; apenas a moldura empírica
das recorrências era excessiva. & \feito. A antiga comparação de seletores virou
regressão interna; foram retiradas inferências de desempenho e superioridade; o bloco de
PD foi apresentado como teoria auxiliar de exatidão.\\
M-5 & O modelo unitário precisava de cota de crescimento dos racionais e custo de
Turing. & \feito. Definidos comprimento binário racional, $B_{\mathcal Z}$,
$B_{\rm cost}$ e $\mathsf M_{\mathbb Z}$; dadas cotas de tempo e memória em bits para as
duas recorrências.\\
M-6 & A regra de empate altera o valor exato $49/64$. & \feito. O enunciado agora escolhe
o menor ponto da grade, equivalentemente o menor numerador no denominador comum.\\
B-1 & Uma falsa edição IEEE em português não podia permanecer. & \feito. Mestre, PDF e
estilos locais IEEE sem uso removidos. Há um único PDF nacional/unificado verdadeiro. Nenhuma tradução
foi simulada.\\
B-2 & O documento-sistema deve sobreviver, mas a publicação deve seguir A/B/C/D e R. &
\parcial. O documento integral foi preservado e os mapas A/B/C/R foram congelados. A
extração final, tradução de B/C, escolha de veículo e o DLPG são etapas externas ao
pacote-base.\\
B-3 & A introdução precisava de 3--4 páginas e posicionamento reconhecível. & \feito.
Introdução reescrita com problema, vizinhança bibliográfica, não-reivindicações, oito
contribuições demonstradas e organização editorial.\\
E-1 & Doze floats precisavam de chamada explícita. & \feito. As sete figuras e cinco
tabelas listadas no parecer passaram a ser chamadas no corpo; o auditor detecta floats
ativos sem referência.\\
E-2 & Remover vestígios de um primeiro artigo e renomear o rótulo histórico. & \feito.
Expressão removida; \texttt{tab:part9-audit} virou
\texttt{tab:dyadic-constants-audit}.\\
E-3 & Os artefatos de \texttt{restricoes\_rede.py} estavam órfãos. & \feito. A tabela da
tenda iterada e a figura de restrições de rede foram inseridas e interpretadas na Seção
20.\\
E-4 & Sete resultados centrais não eram retomados. & \feito. Foram adicionadas chamadas
narrativas aos teoremas de nós livres, refinamento adaptativo, gargalo, complexidade de
estados, imersão direta, pente crítico e estabilidade arquitetural.\\
E-5 & Resumos deviam ter no máximo 250 palavras e nenhuma equação destacada. & \feito.
Resumo e abstract foram reescritos; contagem simples do fonte: 213 e 202 palavras,
respectivamente; nenhum display.\\
E-6 & Faltavam declarações editoriais e metadados. & \feito no limite da informação
disponível. Agradecimentos, financiamento, conflitos, dados/código, ORCID e nota
biográfica foram adicionados. Campos não informados foram marcados como tais; palavras-
chave do PDF foram preenchidas.\\
E-7 & O texto comparava $49/64$ com $3/(2N^2)$ em normalizações diferentes. & \feito.
Depois da multiplicação por $N^2$, o certificado é escrito como $3/2$; a forma não
normalizada é dada entre parênteses.\\
E-8 & O pacote precisa de depósito público com identificador persistente. & \externo.
O pacote está pronto e a declaração de disponibilidade foi incluída, mas nenhum DOI,
Zenodo ou arXiv foi inventado.\\
J-1 & $\Gamma$ define semântica operacional; $\Gamma^\sharp$ é codec. & \feito. Novo
parágrafo descreve as duas conversões e prova a preservação da célula e da função.\\
J-2 & A inclinação não nula sustenta a recuperação de célula. & \feito. A Seção 3 agora
explícita a necessidade; $\LibraryMarked$ é declarada apenas uma extensão auxiliar.\\
J-3 & O refinamento adaptativo tinha relevância externa, mas precisava de amarração
interna. & \feito. O resultado é retomado na introdução, no corpo e na conclusão.\\
P-1--P-3 & Dependência, sobreposição e título do DLPG dizem respeito a outro manuscrito.
& \externo. Os mapas registram a ordem B antes de D e o título congelado de B. Não houve
edição silenciosa de um pacote DLPG que não integra o ZIP-base.\\
G-1--G-8 & Guardas contra enfraquecimento ou falsa conclusão editorial. & \feito no
documento-sistema: nenhum teorema removido; Seção 8.6 mantida; $\Gamma$ e
$\Gamma^\sharp$ preservados; M-3 fecha em $\Library$; M-1 foi reformulado, não apenas
citado; unificado mantido; itens externos registrados como pendentes.\\
\bottomrule
\end{longtable}
\normalsize

\section{Correções matemáticas em detalhe}

\subsection{M-1: funcional de realização e escopo da recíproca}

\textbf{Problema anterior.} A formulação permitia ler a equivalência entre o certificado
e Besov como uma propriedade característica da biblioteca. No entanto, a cota inferior
decorre apenas da desigualdade triangular para segundas diferenças.

\textbf{Substituição conceitual.} A Seção 20 agora introduz explicitamente o objeto como
funcional de realização, cita Hristov e Ivanov (1990) e separa:
\[
 \frac14\,\omega_2(f,N^{-1})_\infty
 \leq \mathcal R_N^{\Library,H}(f)
\]
-- universal para qualquer família de candidatos -- do majorante de Jackson, que exige a
construção racional da biblioteca e sua agenda de altura.

\textbf{Alterações concretas.}
\begin{itemize}
  \item novo parágrafo de enquadramento no início da Seção 20;
  \item novo Lema ``desigualdade inversa universal'';
  \item nova Observação ``o que o funcional mede'';
  \item teorema bilateral renomeado e redigido com o lado direto específico em destaque;
  \item resumo, introdução e conclusão alterados para impedir a reivindicação de
  exclusividade;
  \item referência Hristov--Ivanov incorporada e citada no ponto de uso.
\end{itemize}

\subsection{M-2: uniformidade sobre bolas de Besov}

\textbf{Problema anterior.} O Lema de Jackson possuía um limiar
$n_0(f,K,s)$. Isso sustentava inclusões de conjuntos, mas não uma estimativa uniforme em
um conjunto normado.

\textbf{Nova construção.} Para os nós $t_j$, foram introduzidas grades alternadas
\[
 G_j=\frac1Q\left(\mathbb Z+\frac{j\bmod2}{2}\right),
\]
com escolha do ponto mais próximo e desempate para o menor. Grades consecutivas são
disjuntas, portanto os valores nodais consecutivos nunca coincidem e todas as inclinações
locais permanecem não nulas. A janela
\[
 \frac{H}{16(R+1)}<Q\leq\frac{H}{8(R+1)}
\]
controla simultaneamente erro de quantização, valores centrais e denominadores.

A dependência das primeiras diferenças é rastreada pela desigualdade de Marchaud nos três
regimes $0<s<1$, $s=1$ e $1<s<2$. Disso resulta uma constante
$A_{K,s}$ independente de $f$ tal que, na bola
$\|f\|_{B^s_{\infty,\infty}}^*\leq B$,
\[
 H\geq A_{K,s}(B+1)n^{\max\{1,s\}}
\]
elimina o limiar individual e fornece erro de quantização no máximo
$8(B+1)/H$.

\textbf{Resultado novo.} O Corolário 19.3 define uma única agenda
$\widehat H_{s,\kappa}^{[B]}$ para cada bola e fornece cotas uniformes dos
funcionais de taxa para o endpoint, o regime subcrítico e o regime crítico com
$\kappa>1/q$. A prova trata separadamente o termo do módulo e o termo de
quantização, inclusive a escala $N=1$.

\subsection{M-3: agudeza dentro da biblioteca original}

\textbf{Problema anterior.} O pente de tendas tinha regiões planas e, portanto, pertencia
apenas a $\LibraryMarked$, enquanto a conclusão era anunciada para $\Library$.

\textbf{Perturbação implementada.} Para a aproximação parcial $f_j$, cada folha plana
$J$ da árvore em pente é subdividida uma vez. Sobre $J$ é adicionada uma tenda diádica
nula nos extremos e de altura
\[
 \delta_j=2^{-a_{j+1}-2}.
\]
Logo,
\[
 g_j=f_j+\sum_{J\ \mathrm{plano}}\psi_{J,j}
\]
é contínuo e possui inclinação não nula em todas as folhas refinadas. Portanto
$g_j\in\Library$, e não apenas em $\LibraryMarked$.

\textbf{Orçamentos verificados.}
\begin{itemize}
  \item o número de folhas contínua $O(r_j+j)$;
  \item a nova inclinação em profundidade $d$ tem módulo $2^{d+1}\delta_j$;
  \item a altura é limitada por
  $\max\{2^{a_{j+1}+1},2^{r_j-a_{j+1}}\}$;
  \item o erro adicional uniforme é exatamente controlado por $\delta_j$;
  \item com $N_j=2^{J_0j}$, altura $O(N_j^\gamma)$ e erro
  $O(N_j^{-\alpha})$ são preservados;
  \item as segundas diferenças nas meias larguras das tendas mantém a exclusão de toda
  regularidade acima de $\alpha/(\alpha+\gamma)$ e do endpoint com $q<\infty$.
\end{itemize}

O programa \path{reproducibility/selecao/reversibilidade_regularizada.py} foi
alterado para construir esses átomos da biblioteca original, conferir inclinações não
nulas e registrar folhas, altura e erro em aritmética exata.

\subsection{M-5: crescimento de operandos e complexidade de bits}

A nova Subseção 14.2 define, para $r=p/q$ reduzido,
\[
 \operatorname{bl}(r)=1+\lceil\log_2(|p|+1)\rceil
 +\lceil\log_2(q+1)\rceil.
\]
Se $B_{\mathcal Z}$ limita os dados racionais e $L$ limita o fluxo recursivo, o
novo teorema demonstra a cota conservadora
\[
 B_{\rm cost}=O\!\left(
 P(B_{\mathcal Z}+D+|k_0|+\log(2H))
 +B_{\mathcal Z}+D+\log(2L)\right).
\]
O argumento acompanha avaliação afim, subtração, quadrado, peso, soma de até $P$
frações e penalidades. Em seguida multiplica o custo de uma operação racional e de uma
comparação shortlex pelo número de transições do algoritmo. Foram dadas cotas explícitas
de tempo e memória em bits, com e sem orçamento de folhas, e é registrada a natureza
pseudo-polinomial em $H$ e $2^D$.

\subsection{M-4 e M-6: alcance algorítmico e determinismo}

O bloco de seleção não foi descartado porque o plano preserva o catálogo finito, o
desempate e a totalidade da seleção. O que mudou foi a inferência: ``comparação dos quatro
seletores'' foi substituída por ``instância de regressão e rotas de controle''. Tabela e
figura agora dizem expressamente que não são benchmark, medição de tempo, estudo de escala
ou demonstração de superioridade.

O desempate do Teorema do orçamento conjunto deixou de ser ``fixado'' sem especificação.
Agora, em caso de equidistância, escolhe-se o menor ponto da grade -- ou o menor numerador
no denominador comum. Essa é exatamente a regra necessária para reproduzir $49/64$.

\section{Correções editoriais e estruturais}

\subsection{Substituições textuais verificáveis}

\begin{longtable}{L{0.39\textwidth}L{0.49\textwidth}}
\toprule
\textbf{Antes} & \textbf{Depois}\\
\midrule
\endfirsthead
\toprule
\textbf{Antes} & \textbf{Depois}\\
\midrule
\endhead
``a estimativa de Whitney do primeiro artigo'' & Referência interna direta à estimativa e
à equação correspondentes no documento unificado.\\
\texttt{tab:part9-audit} & \texttt{tab:dyadic-constants-audit}.\\
``com um desempate fixado'' & ``no empate, escolha o ponto menor'', com equivalência ao
menor numerador explicitada no enunciado.\\
``$N^2\|f_\star-g_N\|_\infty=49/64$, enquanto o certificado é $3/(2N^2)$'' &
Comparação em uma única normalização: $49/64$ versus $3/2$, seguida da forma não
normalizada $3/(2N^2)$.\\
``Comparação dos quatro seletores'' & ``Instância de regressão e rotas de controle'', sem
generalização empírica.\\
Requisito Python 3.12+ & Python 3.10 ou posterior, preservando o registro do ambiente de
referência 3.12.13.\\
Dois mestres, inclusive um IEEE em português & Um único mestre ativo e um único PDF
unificado; material IEEE enganoso removido.\\
\bottomrule
\end{longtable}

\subsection{Introdução, resumos e conclusão}

A introdução foi ampliada de uma abertura curta para uma estrutura editorial completa:
problema, vizinhança bibliográfica, fronteira entre resultados clássicos e novos,
não-reivindicações, oito contribuições demonstradas e organização do documento-sistema.
Foram incorporadas referências a splines de nós livres, redes com pesos codificáveis,
códigos universais, segmentação dinâmica e funcionais de realização.

Os dois resumos foram reduzidos abaixo de 250 palavras, sem equações destacadas. Eles
agora mencionam a uniformidade sobre bolas, a agudeza na biblioteca original e a
universalidade da direção inversa. A conclusão foi reescrita para repetir essas
qualificações, preservar a arquitetura A/B/C/R e retirar qualquer leitura de benchmark
da camada de PD.

\subsection{Floats, resultados e artefatos}

Passaram a ter chamadas explícitas as figuras de estado de fronteira, envelope $C^2$,
agendas de altura, fases por partes, dente de serra quadrático, reversibilidade
regularizada e regressão dos seletores; e as tabelas de reversão aproximada, auditoria de
PD, certificado da quina, auditoria de constantes e caso quadrático exato.

Os artefatos de \texttt{restricoes\_rede.py} deixaram de ser órfãos: a tabela da tenda
iterada e a figura de crescimento de recursos foram inseridas na análise da volta de redes
profundas. Sete resultados centrais passaram a ser retomados após o enunciado. As
auditorias estáticas agora falham se esses vínculos desaparecerem.

\subsection{Declarações e metadados}

Foi criado \path{artigo/declaracoes.tex}, contendo agradecimentos, financiamento,
conflitos, disponibilidade de dados/código, ORCID e nota biográfica. Como o material
fornecido não continha confirmações para financiamento, conflito ou ORCID, o texto
registra ``informação não fornecida'' e exige confirmação antes da submissão. Isso evita
atribuir declarações ou identificadores ao autor sem evidência.

O PDF final possui título, autor, assunto e palavras-chave. O bloco de disponibilidade
aponta para os scripts, JSON, tabelas e instruções dentro do próprio pacote, e deixa DOI
como pendência posterior ao depósito público.

\section{Arquitetura editorial resultante}

O antigo plano de três artigos foi substituído pela arquitetura revisada sem destruir o
documento-sistema:

\begin{longtable}{L{0.08\textwidth}L{0.29\textwidth}L{0.50\textwidth}}
\toprule
\textbf{Peça} & \textbf{Escopo congelado} & \textbf{Estado no pacote}\\
\midrule
\endfirsthead
\toprule
\textbf{Peça} & \textbf{Escopo congelado} & \textbf{Estado no pacote}\\
\midrule
\endhead
A & Partes de 5--6 e seções 16--17: constantes e fases diádicas. & Mapa autônomo em
\path{publicacao/A_constantes_diadicas.md}; núcleo matemático corrigido e auditado no
documento-sistema.\\
B & Seções 2--5, 7--8 e verificação pertinente: biblioteca aritmética, endereçamento e
codificação. & Título e ordem congelados em
\path{publicacao/B_biblioteca_aritmetica.md}; J-1/J-2 e referências já integrados.\\
C & Seção 6 e seções 18--20: estratos, inversos e realização. & Mapa em
\path{publicacao/C_estratos_altura.md}; M-1, M-2 e M-3 concluídos no texto integral.\\
D & DLPG, dependente de B. & Não faz parte do ZIP-base corrigido. A ordem de dependência é
registrada; identificador de B não foi inventado.\\
R & Seções 9--15: seleção exata por PD. & Mapa em
\path{publicacao/R_selecao_exata.md}; moldura de exatidão aplicada no documento-sistema.\\
\bottomrule
\end{longtable}

Os mapas não são apresentados como artigos já submetidos. Para publicação internacional,
B e C ainda exigem extração autônoma, tradução integral, ajuste ao limite do veículo e
revisão final do autor. Essa separação entre conteúdo corrigido e ato editorial futuro é
deliberada.

\section{Fortalecimento bibliográfico}

Foram acrescentadas dez entradas e feitas citações no ponto de uso:

\begin{longtable}{L{0.21\textwidth}L{0.25\textwidth}L{0.42\textwidth}}
\toprule
\textbf{Referência} & \textbf{Chave} & \textbf{Função editorial/matemática}\\
\midrule
\endfirsthead
\toprule
\textbf{Referência} & \textbf{Chave} & \textbf{Função editorial/matemática}\\
\midrule
\endhead
Hristov--Ivanov (1990) & \path{hristovivanov1990} & Realização de funcionais em
aproximação restrita; indispensável a M-1.\\
Guhring--Raslan (2021) & \path{guhringraslan2021} & Pesos codificáveis e comparação
com agendas de altura.\\
Burchard (1974) & \path{burchard1974} & Constante clássica e vantagem de nós livres.\\
de Boor (1973) & \path{deboor1973variable} & Aproximação por splines com nós
variáveis.\\
Elbrachter et al. (2021) & \path{elbrachter2021} & Teoria de aproximação por redes
profundas.\\
Petersen--Voigtlaender (2018) & \path{petersenvoigtlaender2018} & Aproximação ReLU de
funções por partes suaves.\\
Elias (1975) & \path{elias1975} & Atribuição dos códigos universais de inteiros.\\
Cover--Thomas (2006) & \path{coverthomas2006} & Fundamentos de informação e Kraft.\\
Perez--Vidal (1994) & \path{perezvidal1994} & Segmentação poligonal por programação
dinâmica posterior a Bellman/Stone.\\
Kessebohmer--Niemann (2025) & \path{kessebohmerniemann2025} & Ordem assintótica de
aproximação adaptativa generalizada.\\
\bottomrule
\end{longtable}

O arquivo bibliográfico passou de 36 para 46 chaves citadas no documento ativo. A
integração foi validada por BibTeX e por busca das referências no texto extraído do PDF.

\section{Livro-razão de arquivos alterados}

\begin{longtable}{L{0.38\textwidth}L{0.51\textwidth}}
\toprule
\textbf{Arquivo ou grupo} & \textbf{Alteração}\\
\midrule
\endfirsthead
\toprule
\textbf{Arquivo ou grupo} & \textbf{Alteração}\\
\midrule
\endhead
\path{artigo/artigo_nacional.tex} & Mestre único; abertura do residual R e declarações
editoriais inseridas.\\
\path{artigo/artigo_ieee.tex} e \path{artigo/estilo_ieee/} & Removidos: o corpo em
português e os estilos sem mestre ativo não podiam ser apresentados como edição IEEE.\\
\path{artigo/resumo_pt.tex}, \path{artigo/abstract_en.tex} & Resumos reduzidos e
reposicionados.\\
\path{artigo/declaracoes.tex} & Novo bloco de conformidade autoral/editorial.\\
\path{artigo/secoes/01_introducao.tex} & Reescrita substancial com estado da arte,
não-reivindicações, contribuições e arquitetura.\\
\path{artigo/secoes/03_gerador.tex} & Papel essencial da inclinação não nula e uso apenas
auxiliar de $\LibraryMarked$.\\
\path{artigo/secoes/06_arvores_besov.tex},
\path{07_adaptativo.tex}, \path{16_constante_diadica.tex} & Novas citações e amarrações
internas aos resultados.\\
\path{artigo/secoes/08_complexidade.tex} & J-1, citações de codificação e regra exata de
desempate.\\
\path{artigo/secoes/11_programacao_arvore_fixa.tex},
\path{12_programacao_arvore_variavel.tex} & Chamadas de figura/teorema e literatura de
segmentação.\\
\path{artigo/secoes/14_complexidade_estabilidade.tex} & Teorema e prova de complexidade
de bits (M-5).\\
\path{artigo/secoes/15_validacao_selecao.tex} & Reclassificação de comparação para
regressão exata (M-4), chamadas de tabela/figura/teorema.\\
\path{artigo/secoes/17_casos_exatos.tex} & Chamadas de floats e renomeação do rótulo
histórico.\\
\path{artigo/secoes/19_teoremas_diretos.tex} & Versão uniforme do Jackson, agenda de bola
e imersões quantitativas (M-2); correção do vestígio textual.\\
\path{artigo/secoes/20_reversibilidade.tex} & M-1 e M-3; inserção dos artefatos de rede;
novas chamadas narrativas.\\
\path{artigo/secoes/21_verificacao_computacional.tex} & Python 3.10+, normalização E-7 e
novos testes do pente original.\\
\path{artigo/secoes/22_conclusao.tex} & Conclusão coerente com a nova moldura e a
arquitetura editorial.\\
\path{artigo/secoes/R_abertura_pd.tex} & Novo aviso de escopo antes do bloco de PD.\\
\path{references.bib} & Dez referências novas.\\
\path{reproducibility/selecao/reversibilidade_regularizada.py} & Auditoria executável da
perturbação de M-3 e novos campos JSON.\\
\path{auditoria/auditar_fontes.py} & Novas proibições, exigências de chamadas, chaves
bibliográficas e verificações M-1--M-6/J-1--J-3.\\
\path{auditoria/auditar_pdfs.py} & Verificação do único PDF, metadados, termos obrigatórios
e referências incorporadas.\\
\path{compilar.sh}, \path{verificar.sh} & Pipeline de um único documento, execução das
seis suítes, auditoria do PDF e manifesto.\\
\path{publicacao/*.md} & Novos mapas A/B/C/R e guarda-corpos.\\
\path{README.md} & Instrução do pacote corrigido, Python 3.10+, arquitetura e pendências
externas.\\
Figuras, tabelas e JSON gerados & Regenerados pelas suítes; os artefatos de restrições de
rede passaram a integrar o PDF.\\
\bottomrule
\end{longtable}

\section{Validação executada}

\subsection{Provas automatizadas e aritmética exata}

As seis suítes do pacote foram executadas com código de saída zero:
\begin{enumerate}
  \item \texttt{reproduzir\_curvas.py};
  \item \texttt{seletor\_dinamico\_exato.py};
  \item \texttt{constante\_diadica.py};
  \item \texttt{espacos\_aproximacao.py};
  \item \texttt{restricoes\_rede.py};
  \item \texttt{reversibilidade\_regularizada.py}.
\end{enumerate}

Na suíte de reversibilidade foram registrados 2.952 round-trips exatos, 84 redes
quantizadas, 1.260 certificados triangulares e seis linhas do pente crítico; todas as
asserções passaram. As demais suítes confirmaram as identidades de programação dinâmica,
constantes quadráticas, envelopes, agendas diretas e restrições de redes.

\subsection{Auditoria estática}

A auditoria final reportou:
\begin{center}
\fbox{\parbox{0.88\textwidth}{29 fontes ativos; 413 rótulos; 281 referências internas;
46 chaves citadas; 97 ambientes matemáticos.}}
\end{center}

Ela também procura vestígios proibidos, chaves não resolvidas, floats não chamados,
rótulos duplicados, referências bibliográficas ausentes e as formulações obrigatórias
das correções.

\subsection{Compilação e inspeção do PDF}

O documento foi compilado por \texttt{latexmk}/pdfLaTeX/BibTeX até estabilização das
referências. Não restaram citações ou referências indefinidas, rótulos múltiplos,
\texttt{Overfull \\hbox} ou \texttt{Overfull \\vbox}. O PDF final tem:
\begin{itemize}
  \item 77 páginas A4;
  \item 868.856 bytes na compilação auditada;
  \item SHA-256
  \nolinkurl{81ae8a5e50e12222d85f95378436b49facbefa67e14b631f8df364af8d9edf9a};
  \item título, autor, assunto e palavras-chave preenchidos;
  \item nenhum formulário, JavaScript, criptografia ou rotação de página.
\end{itemize}

Todas as 77 páginas foram renderizadas em PNG e examinadas em folhas de contato. Também
foram ampliadas as páginas da complexidade de bits, do Jackson uniforme, das imersões
sobre bolas, da volta ReLU, do pente perturbado, do funcional de realização, das
declarações e da bibliografia. Não foram observados cortes, sobreposições, glifos
corrompidos, tabelas ilegíveis ou figuras sem legenda.

\section{Pendências que não podiam ser preenchidas com rigor}

Os itens seguintes permanecem deliberadamente abertos:

\begin{itemize}
  \item \textbf{DOI/Zenodo e arXiv.} Exigem depósito em serviços externos e decisão do
  autor. Nenhum identificador foi criado ou antecipado no texto.
  \item \textbf{ORCID.} Não foi fornecido; nenhum número foi inferido.
  \item \textbf{Financiamento e conflitos.} Requerem declaração expressa do autor; o
  pacote apenas marca a informação como não fornecida.
  \item \textbf{Biografia completa.} Somente os fatos já presentes no manuscrito foram
  mantidos.
  \item \textbf{Tradução integral de B e C.} O antigo arquivo IEEE em português foi
  removido. Uma edição internacional só deve ser criada quando todo o corpo, figuras,
  tabelas e metadados forem traduzidos e revisados.
  \item \textbf{Extração dos manuscritos A/B/C e edição do DLPG.} Os escopos e
  dependências foram congelados, mas a extração editorial final depende de veículo,
  limite de páginas, idioma e revisão autoral. O DLPG não integra o pacote-base.
  \item \textbf{Alegações empíricas.} Um estudo de escala, tempos, dados reais e
  comparação ampla poderia sustentar um futuro artigo de métodos; não foi fabricado a
  partir das pequenas instâncias existentes.
\end{itemize}

Essas pendências não afetam as correções matemáticas M-1, M-2, M-3, M-5 e M-6 nem a
reclassificação honesta de M-4. Elas impedem apenas que o pacote seja descrito como
submissão internacional já finalizada.

\section{Conteúdo do ZIP corrigido}

O arquivo final contém:
\begin{itemize}
  \item o PDF unificado corrigido;
  \item todos os fontes LaTeX ativos;
  \item bibliografia;
  \item scripts, JSON, figuras e tabelas de reprodutibilidade;
  \item auditorias de fontes e PDF;
  \item scripts de compilação e verificação;
  \item mapas editoriais A/B/C/R;
  \item este relatório em PDF e em fonte LaTeX;
  \item manifesto SHA-256 recalculado sobre os arquivos distribuídos.
\end{itemize}

O comando \path{./verificar.sh} verifica o manifesto, reexecuta a auditoria estática e
inspeciona o PDF principal. O comando \path{./compilar.sh} regenera figuras e tabelas,
executa as seis suítes, recompila o artigo, audita o resultado e atualiza o manifesto.

\section{Conclusão de auditoria}

As orientações foram cruzadas sem tratar o plano revisado como cancelamento indevido do
parecer. Os achados matemáticos foram mantidos; M-3 foi reforçado; M-4 foi limitado ao
que de fato estava frágil; e a arquitetura editorial foi atualizada para A/B/C/D com R,
preservando o documento-sistema. O pacote entregue contém uma versão matematicamente
fortalecida, editorialmente mais honesta é automaticamente verificável. As únicas
pendências declaradas são aquelas que exigem ato externo, informação autoral ou uma
decisão de submissão que não pode ser inferida dos arquivos.

\end{document}
